\chapter*{Notation}
\label{c:Notation}

Note that, to be consistent with the Crohn's disease case 
study used throughout this thesis, the endpoint is oriented 
so that lower values indicate a better clinical outcome.


\section*{Parameters and data}

\begin{tabularx}{\linewidth}{@{}lX@{}}
\toprule
Symbol & Meaning \\
\midrule
$\theta_C,\ \theta_T$        & true control and treatment means \\
$\bar y_C,\ \bar y_T$        & observed arm means \\
$\delta$                     & treatment contrast, $\delta = \theta_C - \theta_T$ (lower $y$ is beneficial) \\
$\sigma_C,\ \sigma_T$        & per patient standard deviations, assumed known \\
$\sigma$                     & their common value, where the two arms are taken equal \\
$n_C,\ n_T$                  & control and treatment sample sizes \\
$n_{\text{trial}} = n_C + n_T$ & total sample size of a fixed design \\
$n_A$                        & effective sample size of the historical information \\
$v_C^2 = \sigma_C^2/n_C$       & sampling variance of $\bar y_C$ \\
$v_T^2 = \sigma_T^2/n_T$       & sampling variance of $\bar y_T$ \\
$r_T = n_T/n_C$              & allocation ratio \\
$r_A = n_A/n_C$              & historical to control effective sample size ratio \\
\bottomrule
\end{tabularx}

\section*{Priors}

\begin{tabularx}{\linewidth}{@{}lX@{}}
\toprule
Symbol & Meaning \\
\midrule
$\pi_A(\theta_C)$                       & analysis prior, used to form the posterior \\
$\pi_D(\theta_C)$                       & design prior, used to average operating characteristics \\
$\mu_{A,C},\ \sigma_{A,C}^2$            & mean and variance of the analysis prior \\
$\mu_{D,C},\ \sigma_{D,C}^2$            & mean and variance of the design prior \\
$W = \sigma_{A,C}^2/(\sigma_{A,C}^2+v_C^2)$ & weight the data carries in the control estimate \\
$1-W$                                   & weight the historical prior carries \\
$\omega$                                & weight of the robust component in a robust mixture prior \\
$\mu_R,\ \sigma_R^2$                    & mean and variance of that robust component, a UIP with effective sample size one \\
\bottomrule
\end{tabularx}

\section*{Posterior and decision quantities}

\begin{tabularx}{\linewidth}{@{}lX@{}}
\toprule
Symbol & Meaning \\
\midrule
$\mu_{\text{post},C},\ \sigma_{\text{post},C}^2$ & mean and variance of the control posterior \\
$\hat\delta$                                 & posterior mean of the contrast \\
$V_\delta$                                   & posterior variance of the contrast \\
$V_{\text{samp}}$                            & sampling variance of $\hat\delta$ at a fixed $\theta_C$ \\
$V_{\text{marg}}$                            & marginal variance of $\hat\delta$ under the design prior \\
$p$                                          & posterior probability threshold for success \\
$\Phi(\cdot),\ \phi(\cdot)$                  & standard Normal distribution and density functions \\
$z_p = \Phi^{-1}(p)$                         & corresponding standard Normal quantile \\
$q$                                          & null hypothesis threshold for $\delta$, e.g.\ $q=0$ for a superiority test \\
$\DV$                                        & decision value of the dual criterion \\
$k = q + z_p\sqrt{V_\delta}$                 & decision threshold in the Normal case: success iff $\hat\delta > k$ \\
$K$                                          & total number of analyses, interims plus final \\
$t_k$                                        & information fraction at look $k$ \\
\bottomrule
\end{tabularx}

\section*{Sequential designs and calibration}

\begin{tabularx}{\linewidth}{@{}lX@{}}
\toprule
Symbol & Meaning \\
\midrule
$\mathcal{S} = \{(n_{T,k}, n_{C,k})\}_{k=1}^{K}$ & schedule, the cumulative arm sizes at each of the $K$ looks \\
$n_{m,k}$                    & cumulative size of arm $m \in \{T,C\}$ at look $k$ \\
$n_{\max}$                   & maximum total sample size, reached if the trial runs to the final look \\
$N$                          & realized total sample size, a random variable in a sequential design \\
$N^+,\ N^-$                  & $N$ conditional on success and on no success \\
$p_{k,\text{eff}},\ q_{k,\text{eff}}$ & efficacy threshold and reference value at look $k$ \\
$p_{k,\text{fut}},\ q_{k,\text{fut}}$ & futility threshold and reference value at look $k$ \\
$p_{\mathrm{HP}}$            & near impossible interim threshold of the Haybittle Peto boundary \\
$\Delta$                     & shape parameter of the Wang Tsiatis boundary family \\
$\delta_{\text{target}}$     & effect size a design is calibrated to, the target for avgPower \\
$\alpha^\ast,\ 1-\beta^\ast$ & the avgT1E and avgPower a design is calibrated to \\
\bottomrule
\end{tabularx}

\section*{Operating characteristics and metrics}

\begin{tabularx}{\linewidth}{@{}lX@{}}
\toprule
Symbol & Meaning \\
\midrule
$\alpha$                & nominal significance level \\
$\TIE(\theta_C)$        & conditional type I error at a fixed $\theta_C$ \\
$\avgTIE$               & average type I error, integrated over $\pi_D$ \\
$\Pow(\theta_C,\delta)$ & conditional power \\
$\avgPow(\delta)$       & average power at effect $\delta$, integrated over $\pi_D$ \\
$\DOR$                  & diagnostic odds ratio \\
$\LRpos,\ \LRneg$       & likelihood ratios of a positive and a negative outcome \\
$\EUII$                 & experimental unit information index, $\DOR^{1/n_{\text{trial}}}$ for a fixed design and $(\LRpos)^{\Erw[1/N^+]} / (\LRneg)^{\Erw[1/N^-]}$ for a sequential one \\
$\Lodd(x) = \log\Phi(x)-\log\Phi(-x)$ & log odds of $\Phi(x)$, used in the asymptotic EUII derivation \\
$\Mills(x) = (1-\Phi(x))/\phi(x)$ & Mills ratio, used in the same derivation \\
$\Pr(H_1)$              & prior probability of the alternative. \\
\bottomrule
\end{tabularx}

\section*{Abbreviations}

\begin{tabularx}{\linewidth}{@{}lX@{}}
\toprule
Abbreviation & Meaning \\
\midrule
MAP                     & meta analytic predictive \\
UIP                     & unit information prior \\
ESS                     & effective sample size \\
$\SC$                   & single criterion \\
$\DC$                   & dual criterion \\
\bottomrule
\end{tabularx}

%%% Local Variables: 
%%% mode: latex
%%% TeX-master: "MasterThesisSfS"
%%% End: 
